\section{Experimental setup}
\label{sec:setup}

\subsection{Common synthetic protocol}

All synthetic experiments use corrected budget accounting $\Stotal=\Rsettings \Nshots T$ with $\Rsettings=2$ measurement settings per labeled time step. The base reservoir uses $Q=4$ qubits, depth $1$, a ring-CNOT entangler, leaky integration with $\eta=0.2$, window length $L=10$, chronological train/validation/test splits with washout $t_0=30$, and one-step-ahead forecasting on normalized benchmarks. Unless stated otherwise, the confirmation set consists of $20$ seeds.

The common synthetic sweep varies
\[
N_{\mathrm{shots}}\in\{10,12,15,18,20,25,30,40,50\}
\]
with no additional synthetic noise on top of shot sampling. The synthetic suite fixes the ridge parameter globally at $\lambda=10$, drawn from the same common log-grid used throughout the study. This sweep is used to probe how the EV-minus-Split gap changes along the row-limited to well-conditioned regime axis.

\subsection{A single common operating point}

The representative panels in Figure~\ref{fig:rep} and the protocol baselines in Figure~\ref{fig:base} are meant to be maximally transparent rather than presentation-driven. We therefore use one common synthetic operating point across all three tasks instead of task-specific settings.

The selection rule is explicit: choose the largest $N_{\mathrm{shots}}$ value from the common sweep for which all three benchmarks simultaneously satisfy two conditions in the $20$-seed mean, namely EV remains below $\NRMSE=1$ and Split improves on EV. Under that rule, the shared operating point is $Q=4$, $N_{\mathrm{shots}}=18$, $\sigma=0$, $\lambda=10$, and $\rhoEV=2.65125$.

This rule removes the appearance of magic numbers. The operating point is not chosen separately for Mackey--Glass, Lorenz, and NARMA10; it is the single largest shared point in the sweep that still keeps the EV baseline in a reasonable forecasting regime on all three tasks. The left panels in Figure~\ref{fig:rep} then show one representative seed per task, chosen as the seed whose Split test $\NRMSE$ is closest to the 20-seed median, while the right panels still report all 20 seed outcomes.

\subsection{Empirical regime plot}

The empirical regime plot (Figure~\ref{fig:gap}) is a one-parameter sweep through the common synthetic setup. Its role is narrower than that of the representative panels: it shows how the gap
\[
\Delta = \NRMSE(\mathrm{EV})-\NRMSE(\mathrm{Split})
\]
changes along a one-parameter slice in $\rhoEV$. Faint points show individual seeds and the darker markers report the $20$-seed mean with a confidence interval. This plot is intended as the direct empirical test of the regime picture from Section~\ref{sec:methodology}.

\subsection{Imported hardware validation}

To address the most obvious reviewer concern, we include a real-device validation from the companion hardware study on \texttt{ibm\_marrakesh} \cite{ahmed2025optimal}. The hardware task is Mackey--Glass with $Q=6$, depth $5$, $L=20$, $\Stotal=20{,}000$, $N_{\mathrm{shots}}=40$, $\lambda=10$, and $\Rsettings=2$. Raw bitstrings were collected for eight reservoir seeds, and EV, Raw, and Split were then evaluated \emph{offline} on the same hardware shot records without error mitigation.

We adopt those hardware results directly rather than re-optimizing them inside the synthetic rewrite. This keeps the hardware section clean: the device evidence is exactly the evidence already obtained on hardware.

\subsection{Larger-reservoir sweep}

The appendix enlarges the reservoir while keeping the protocol fixed. For all three benchmarks we sweep $Q\in\{4,6,8,10\}$ with eight seeds per point; for Mackey--Glass and NARMA10 we additionally include $Q=12$. The purpose is not to claim asymptotic scaling, but to answer the concrete reviewer question of whether the protocol effect survives once the feature dimension becomes larger and the conditioning story becomes more interesting.
